\chapter{Đọc thêm 2: Phân rã ma trận không âm (Ghi chú bài giảng, TU Dortmund)}
\label{ch:M5L1DocThem2}

\section{Thông tin nguồn}

\begin{itemize}
  \item \textbf{Nguồn:} ``Non-negative Matrix Factorization -- Lecture Notes.'' TU Dortmund.
  \item \textbf{Phần cần đọc:} Toàn bộ trang
  \item \textbf{Ghi chú:} bản chuyển ngữ tiếng Việt súc tích, bám cấu trúc nguồn (giữ nguyên tiêu đề, công thức, số liệu, bảng, chú thích và danh mục tài liệu tham khảo).
\end{itemize}

\subsubsection{Phân rã ma trận không âm (NMF, NNMF) {[}LeSe99{]}}

Cho ma trận dữ liệu \(X \in \mathbb{R}^{m \times n}_{\ge 0}\) và tham số \(k\) với \(1 \le k \le \min\{m, n\}\), hãy tìm hai ma trận

\[W \in \mathbb{R}^{m \times k}_{\ge 0}\]

\[H \in \mathbb{R}^{k \times n}_{\ge 0}\]

sao cho

\[X \approx WH\]

\textbf{Biến thể 1: Sai số bình phương} -- cực tiểu hóa

\[\|X - WH\|^2 = \sum_i \sum_j \left(X_{ij} - (WH)_{ij}\right)^2\]

\textbf{Biến thể 2: Độ phân kỳ (divergence)} -- cực tiểu hóa log-hợp lý của \(X_{ij} \sim \text{Poisson}\left((WH)_{ij}\right)\)

\[L(W, H) = -\sum_i \sum_j \left(X_{ij} \log (WH)_{ij} - (WH)_{ij}\right)\]

\subsubsection{NMF: Cập nhật nhân tính cho sai số bình phương}

Phương pháp hạ gradient cho sai số bình phương {[}LeSe00; LeSe99{]}:

\[\|X - WH\|^2 = \sum_i \sum_j \left(X_{ij} - (WH)_{ij}\right)^2\]

\begin{enumerate}
\def\labelenumi{\arabic{enumi}.}
\item
  Khởi tạo bằng các giá trị không âm ngẫu nhiên (về sau: dùng phương pháp heuristic tốt hơn).
\item
  Tối ưu \(W\):
\end{enumerate}

\[w_{ik} \leftarrow w_{ik}\,\frac{(XH^T)_{ik}}{(WHH^T)_{ik}}\]

\begin{enumerate}
\def\labelenumi{\arabic{enumi}.}
\setcounter{enumi}{2}
\item
  Tối ưu \(H\):
\end{enumerate}

\[h_{kj} \leftarrow h_{kj}\,\frac{(W^T X)_{kj}}{(W^T WH)_{kj}}\]

\begin{enumerate}
\def\labelenumi{\arabic{enumi}.}
\setcounter{enumi}{3}
\item
  Lặp lại bước 2-3 cho đến khi mức cải thiện của \(\|X - WH\|^2\) nhỏ hơn một ngưỡng.
\end{enumerate}

Thuật toán có thể mắc kẹt ở điểm bất động cục bộ hoặc điểm yên ngựa, nên không đảm bảo tìm được nghiệm tối ưu. Lưu ý: cần \(w_{ij} > 0\), \(h_{ij} > 0\), vì vậy thường cộng một hằng số nhỏ \(\varepsilon = 10^{-16}\) vào mọi giá trị.

\subsection{NMF: cập nhật nhân tính cho độ phân kỳ KL}

Phương pháp hạ gradient cho độ phân kỳ KL (KL-Divergence) {[}LeSe00; LeSe99{]}:

\[L(W,H) = \sum_i \sum_j \left( x_{ij} \log (WH)_{ij} - (WH)_{ij} \right)\]

\begin{enumerate}
\def\labelenumi{\arabic{enumi}.}
\item
  Khởi tạo ngẫu nhiên (hoặc bằng heuristic tốt hơn).
\item
  Tối ưu \(W\):
\end{enumerate}

\[w_{ik} \leftarrow w_{ik}\,\frac{\sum_{j=1}^{m} h_{kj}\,\dfrac{x_{ij}}{(WH)_{ij}}}{\sum_{j=1}^{m} h_{kj}}\]

\begin{enumerate}
\def\labelenumi{\arabic{enumi}.}
\setcounter{enumi}{2}
\item
  Tối ưu \(H\):
\end{enumerate}

\[h_{kj} \leftarrow h_{kj}\,\frac{\sum_{i=1}^{n} w_{ik}\,\dfrac{x_{ij}}{(WH)_{ij}}}{\sum_{i=1}^{n} w_{ik}}\]

\begin{enumerate}
\def\labelenumi{\arabic{enumi}.}
\setcounter{enumi}{3}
\item
  Lặp lại bước 2-3 cho đến khi mức cải thiện của \(L(W,H)\) nhỏ hơn một ngưỡng.
\end{enumerate}

Tương tự, thuật toán có thể mắc kẹt ở điểm bất động cục bộ hoặc điểm yên ngựa, không đảm bảo tìm được nghiệm tối ưu. Lưu ý: cần \(w_{ij} > 0\), \(h_{ij} > 0\), nên thường cộng hằng số nhỏ \(\varepsilon = 10^{-16}\) vào mọi giá trị.

\subsection{Các thuật toán khác cho NMF}

Khi cài đặt nguyên văn, các thuật toán trên chạy chậm do mỗi vòng lặp phải nhân ma trận dày đặc với độ phức tạp \(O(n^3)\). Nhiều biến thể đã được phát triển:

\begin{itemize}
\item
  Bình phương tối thiểu luân phiên (Alternating Least Squares) {[}PaTa94{]}
\item
  ALS với gradient chiếu (Projected-gradient ALS) {[}Lin07{]}
\item
  Tối ưu hóa tựa Newton (Quasi-Newton) {[}ZdCi06{]}
\item
  Bình phương tối thiểu luân phiên phân cấp (Hierarchical ALS) {[}CiPh09{]}
\item
  Khởi tạo bằng spherical-\(k\)-means
\item
  Khởi tạo bằng SVD {[}BBLP07{]}
\item
  Khảo sát các phương pháp khởi tạo: {[}LMAC14{]}
\item
  Độ phân kỳ beta (Beta divergence) {[}FéId11{]}
\item
  Phân rã tensor không âm {[}CiPh09{]}
\item
  ...
\end{itemize}

\subsection{pLSI {[}Hofm99{]} là phân rã ma trận không âm {[}GaGo05{]}}

Lập chỉ mục ngữ nghĩa ẩn xác suất (Probabilistic Latent Semantic Indexing, pLSI; xem Phần 4) là phiên bản xác suất của LSI/LSA. PLSI dùng mô hình hỗn hợp dạng đồ thị (ký hiệu mô hình hỗn hợp, khác với Phần 4):

\[P(w_i \mid d_j) = \sum_t P(t)\,P(d_j \mid t)\,P(w_i \mid t)\]

\begin{itemize}
\item
  \(W\): xác suất của từ trong mỗi chủ đề, \(P(w_i \mid t)P(t)\).
\item
  \(H\): xác suất của chủ đề trong mỗi tài liệu, \(P(d_j \mid t)\).
\item
  Các xác suất đều không âm.
\end{itemize}

"Mọi nghiệm cực đại hóa hợp lý (cục bộ) của PLSA đều là nghiệm của NMF với độ phân kỳ KL." {[}GaGo05{]}

\subsection{Lịch sử của phân rã ma trận không âm}

NMF được Daniel D. Lee và H. Sebastian Seung phổ biến qua công trình: Learning the parts of objects by non-negative matrix factorization. In: Nature 401, 4, 1999. DOI: 10.1038/44565.

Các cách tiếp cận tương tự đã xuất hiện trước đó:

\begin{itemize}
\item
  "Phân rã hạng không âm" (Non-negative Rank Factorization)

  \begin{itemize}
    \item
    M.W. Jeter, and W.C. Pye. A note on nonnegative rank factorizations. In: Linear Algebra and its Applications 38, 3, 1981. DOI: 10.1016/0024-3795(81)90018-5
  \item
    Ji-Cheng Chen. The nonnegative rank factorizations of nonnegative matrices. In: Linear Algebra and its Applications 62, 11, 1984. DOI: 10.1016/0024-3795(84)90096-X
  \end{itemize}
\item
  "Phân rã ma trận dương" (Positive Matrix Factorization)

  \begin{itemize}
    \item
    Pentti Paatero, and Unto Tapper. Positive matrix factorization: A non-negative factor
  \end{itemize}
\end{itemize}

model with optimal utilization of error estimates of data values. In: Environmetrics 5, 16, 1994. DOI: 10.1002/env.3170050203

\subsection{Các vấn đề của NMF}

Một phân tích lý thuyết đáng chú ý:

David L. Donoho, and Victoria Stodden. When Does Non-Negative Matrix Factorization Give a Correct Decomposition into Parts? In: NIPS, 8, 2003

Các quan sát:

\begin{itemize}
\item
  Phép phân rã không trực giao (khác với SVD, PCA).
\item
  Phép phân rã không duy nhất (kể cả khi đạt được \(X = WH\); không chỉ do hoán vị và thay đổi tỷ lệ các thành phần).
\item
  Đặc biệt, các vùng bất biến bị phân rã một cách tùy ý (từ dừng \ensuremath{\approx} các thành phần bất biến của dữ liệu).
\item
  Các giá trị 0 gây khó khăn cho chiến lược cập nhật nhân tính (multiplicative update).
\end{itemize}

Cần thêm các ràng buộc bổ sung, chẳng hạn điều chuẩn thưa (sparsity regularization) {[}EgKo04; Hoye04{]}.

\subsection{Tài liệu tham khảo}

{[}BBLP07{]} Berry, M.W., Browne, M., Langville, A.N., Pauca, V.P. and Plemmons, R.J. 2007. Algorithms and applications for approximate nonnegative matrix factorization. Comput. Stat. Data Anal. 52, 1 (2007), 155--173. DOI:10.1016/j.csda.2006.11.006

{[}CiPh09{]} Cichocki, A. and Phan, A.H. 2009. Fast local algorithms for large scale nonnegative matrix and tensor factorizations. IEICE Trans. Fundam. Electron. Commun. Comput. Sci. 92-A, 3 (2009), 708--721. DOI:10.1587/transfun.E92.A.708

{[}EgKo04{]} Eggert, J. and Korner, E. 2004. Sparse coding and NMF. IJCNN (2004), 2529--2533.

{[}FéId11{]} Févotte, C. and Idier, J. 2011. Algorithms for nonnegative matrix factorization with the \(\beta\)-divergence. Neural Computation. 23, 9 (2011), 2421--2456. DOI:10.1162/NECO\_a\_00168

{[}GaGo05{]} Gaussier, É. and Goutte, C. 2005. Relation between PLSA and NMF and implications. SIGIR (2005), 601--602.

{[}Hofm99{]} Hofmann, T. 1999. Learning the similarity of documents: An information-geometric approach to document retrieval and categorization. Neural information processing systems, NIPS (1999), 914--920.

{[}Hoye04{]} Hoyer, P.O. 2004. Non-negative matrix factorization with sparseness constraints. Journal of Machine Learning Research. 5, (2004), 1457--1469.

{[}LeSe00{]} Lee, D.D. and Seung, H.S. 2000. Algorithms for non-negative matrix factorization. NIPS (2000), 556--562.

{[}LeSe99{]} Lee, D.D. and Seung, H.S. 1999. Learning the parts of objects by non-negative matrix factorization. Nature. 401, 6755 (1999), 788--791. DOI:10.1038/44565

{[}Lin07{]} Lin, C.-J. 2007. Projected gradient methods for nonnegative matrix factorization. Neural Computation. 19, 10 (2007), 2756--2779. DOI:10.1162/neco.2007.19.10.2756

{[}LMAC14{]} Langville, A.N., Meyer, C.D., Albright, R., Cox, J. and Duling, D. 2014. Algorithms, initializations, and convergence for the nonnegative matrix factorization. CoRR. abs/1407.7299, (2014).

{[}PaTa94{]} Paatero, P. and Tapper, U. 1994. Positive matrix factorization: A non-negative factor model with optimal utilization of error estimates of data values. Environmetrics. 5, 2 (1994), 111--126. DOI:10.1002/env.3170050203

{[}ZdCi06{]} Zdunek, R. and Cichocki, A. 2006. Non-negative matrix factorization with quasi-newton optimization. Artificial intelligence and soft computing (ICAISC) (2006), 870--879.
